\section{Gambaran Umum Program}

Program lexer untuk bahasa pemrograman Arion didesain sebagai aplikasi berbasis command-line yang menerima masukan berupa file source code dan menghasilkan daftar token sebagai keluaran. Alur umum eksekusi program adalah sebagai berikut.

\begin{enumerate}
    \item \textbf{Validasi Argumen}. Program menerima dua argumen dari command-line, yaitu nama file input (wajib) dan nama file output (opsional). Jika argument tidak sesuai, program akan menampilkan pesan error dan keluar dengan exit code 1.

    \item \textbf{Pembacaan File Input}. Program membuka file input dengan mode baca. Jika file tidak dapat dibuka, program menampilkan error message dan keluar. Seluruh isi file dibaca ke dalam buffer string di memori.

    \item \textbf{Inisialisasi Lexer}. Objek atau struktur Lexer diinisialisasi dengan string source code yang sudah dibaca. State awal DFA disetel ke \texttt{START}, posisi pembacaan diatur ke 0, dan daftar token diinisialisasi kosong.

    \item \textbf{Tokenisasi}. Lexer menjalankan mesin DFA dengan membaca karakter demi karakter dari source code. Untuk setiap karakter, berdasarkan state DFA saat ini, lexer menentukan transisi ke state berikutnya atau menyelesaikan pembacaan satu token dan memancarkannya.

    \item \textbf{Deteksi Error}. Jika lexer menemukan simbol atau urutan karakter yang tidak valid, lexer memancarkan token ERROR dengan pesan diagnostik, namun tetap melanjutkan proses tokenisasi (error recovery).

    \item \textbf{Format Output}. Token-token yang telah dikumpulkan diformat menjadi string output sesuai dengan spesifikasi format lexer. Jika terdapat gap baris pada source code, baris kosong disisipkan pada output untuk mempertahankan struktur visual.

    \item \textbf{Penulisan Output}. Output dicetak ke terminal (standard output). Jika nama file output diberikan sebagai argumen, output juga ditulis ke file tersebut.

    \item \textbf{Kode Keluar}. Program keluar dengan exit code 0 jika tidak ada error leksikal atau exit code 1 jika terdapat error.
\end{enumerate}

Secara arsitektur, program dibagi menjadi beberapa modul yang terpisah untuk meningkatkan modularitas dan maintainability. Modul \texttt{token} mendefinisikan enum TokenType yang berisi 52 jenis token plus ERROR, struktur Token untuk menyimpan informasi token (tipe, nilai, baris, kolom), dan fungsi-fungsi utilitas untuk konversi tipe ke string dan pemformatan output. Modul \texttt{lexer} mengimplementasikan class atau struktur Lexer yang mengenkapsulasi state DFA, buffer pembacaan, dan logika transisi. Modul \texttt{main} bertugas menangani argument parsing, file I/O, dan orchestration keseluruhan proses dari pembacaan input sampai penulisan output.

\section{Algoritma Lexer}

Algoritma lexer pada program ini menggunakan pendekatan \textit{single pass scanning}, yaitu membaca source code sekali dari kiri ke kanan secara karakter demi karakter. Selama proses pembacaan, lexer menyimpan \textit{state} aktif DFA, posisi karakter saat ini, serta buffer lexeme yang sedang dibangun. Setiap karakter yang dibaca akan memicu transisi ke state berikutnya, atau menandakan bahwa sebuah token telah lengkap dan siap dikeluarkan. Dengan pendekatan ini, kompleksitas waktu tokenisasi tetap linear terhadap panjang input, yakni $O(n)$, karena setiap karakter diproses dengan jumlah operasi konstan.

Secara lebih rinci, alur algoritma dimulai dari tahap validasi argumen masukan, kemudian pembacaan file \texttt{input.txt} ke dalam string \texttt{source}. Setelah itu, objek lexer diinisialisasi pada state awal \texttt{START} dengan posisi awal karakter di indeks pertama. Fungsi \texttt{tokenize()} lalu dijalankan dalam loop utama hingga akhir input. Pada setiap iterasi, handler state yang sesuai akan dieksekusi, misalnya handler untuk identifier, angka, string, operator, delimiter, atau komentar. Jika lexeme valid selesai dibaca, lexer memanggil \texttt{emitToken(...)}. Jika ditemukan pola tidak valid, lexer memanggil \texttt{emitError(...)} namun tetap melanjutkan pemrosesan untuk mendukung \textit{error recovery}.

Setelah seluruh input dipindai, program menjalankan \texttt{handleEof()} untuk menangani kasus token yang belum tertutup di akhir file (contohnya komentar atau string yang tidak ditutup). Daftar token yang dihasilkan kemudian diformat menjadi teks keluaran menggunakan fungsi \texttt{formatOutput(...)}. Hasil ini selalu ditulis ke terminal, dan jika argumen file output tersedia, hasil yang sama disimpan ke \texttt{output.txt}. Program mengembalikan exit code 0 jika tidak ada error leksikal, atau 1 jika ditemukan sedikitnya satu token error.

Pseudocode garis besar algoritma program ditunjukkan pada Algoritma~\ref{alg:lexer-main-flow}.

\begin{breakablealgorithm}[Garis Besar Algoritma Lexer (\texttt{input.txt} $\rightarrow$ \texttt{output.txt})]
    \label{alg:lexer-main-flow}
    \begin{algorithmic}[1]
        \Require Argumen command-line: \texttt{argv[1]} = \texttt{input.txt}, \texttt{argv[2]} (opsional) = \texttt{output.txt}
        \Ensure Daftar token tercetak ke terminal dan opsional tersimpan ke \texttt{output.txt}

        \If{jumlah argumen \textless{} 2}
        \State Tampilkan pesan penggunaan program
        \State \Return exit code 1
        \EndIf

        \State Buka berkas \texttt{input.txt}
        \If{berkas gagal dibuka}
        \State Tampilkan pesan error file input
        \State \Return exit code 1
        \EndIf

        \State Baca seluruh isi berkas ke string \texttt{source}
        \State Inisialisasi \texttt{lexer} dengan \texttt{source}
        \State \texttt{tokens} := \texttt{lexer.tokenize()}

        \State \texttt{output} := \texttt{formatOutput(tokens)}
        \State Cetak \texttt{output} ke terminal

        \If{argumen \texttt{output.txt} diberikan}
        \State Buka/buat berkas \texttt{output.txt}
        \If{berkas output gagal dibuka}
        \State Tampilkan pesan error file output
        \State \Return exit code 1
        \EndIf
        \State Tulis \texttt{output} ke \texttt{output.txt}
        \EndIf

        \If{\texttt{lexer.hasErrors()} = true}
        \State \Return exit code 1
        \Else
        \State \Return exit code 0
        \EndIf
    \end{algorithmic}
\end{breakablealgorithm}

Flowchart pada Gambar~\ref{fig:flowchart-lexer-program} memperlihatkan alur eksekusi program secara visual dari penerimaan argumen hingga penentuan exit code.

\begin{figure}[H]
    \centering
    \resizebox{0.85\linewidth}{!}{%
        \begin{tikzpicture}[
                node distance=8mm,
                >=latex,
                line/.style={draw, -latex'},
                term/.style={rectangle, rounded corners, draw, align=center, minimum width=4.4cm, minimum height=8mm},
                proc/.style={rectangle, draw, align=center, minimum width=4.4cm, minimum height=8mm},
                decision/.style={diamond, draw, aspect=2, align=center, inner sep=1pt, minimum width=3.6cm, minimum height=8mm}
            ]

            \node[term] (start) {Mulai Program};
            \node[decision, below=of start] (argc) {Argumen valid?\\(\texttt{argc >= 2})};
            \node[proc, below=of argc] (openin) {Buka \texttt{input.txt}};
            \node[decision, below=of openin] (inok) {File input\\berhasil dibuka?};
            \node[proc, below=of inok] (read) {Baca source code};
            \node[proc, below=of read] (tokenize) {Jalankan \texttt{lexer.tokenize()}};
            \node[proc, below=of tokenize] (format) {Format token ke string output};
            \node[proc, below=of format] (print) {Cetak output ke terminal};
            \node[decision, below=of print] (hasOut) {Ada argumen\\\texttt{output.txt}?};
            \node[proc, below=of hasOut] (writeOut) {Tulis output ke file};
            \node[decision, below=of writeOut] (hasErr) {Ada lexical error?};
            \node[term, below left=12mm and 16mm of hasErr] (end1) {Exit code 1\\(error)};
            \node[term, below right=12mm and 16mm of hasErr] (end0) {Exit code 0\\(sukses)};

            \node[term, right=24mm of argc] (argErr) {Tampilkan usage\\Exit code 1};
            \node[term, right=24mm of inok] (inErr) {Tampilkan error file input\\Exit code 1};

            \draw[line] (start) -- (argc);
            \draw[line] (argc) -- node[right]{Ya} (openin);
            \draw[line] (argc.east) -- ++(6mm,0) node[above,pos=0.35]{Tidak} -- (argErr.west);

            \draw[line] (openin) -- (inok);
            \draw[line] (inok) -- node[right]{Ya} (read);
            \draw[line] (inok.east) -- ++(6mm,0) node[above,pos=0.35]{Tidak} -- (inErr.west);

            \draw[line] (read) -- (tokenize);
            \draw[line] (tokenize) -- (format);
            \draw[line] (format) -- (print);
            \draw[line] (print) -- (hasOut);
            \draw[line] (hasOut) -- node[right]{Ya} (writeOut);
            \draw[line] (hasOut.east) -- ++(8mm,0) |- node[pos=0.25,right]{Tidak} (hasErr.east);
            \draw[line] (writeOut) -- (hasErr);

            \draw[line] (hasErr) -- node[above left]{Ya} (end1);
            \draw[line] (hasErr) -- node[above right]{Tidak} (end0);
        \end{tikzpicture}
    }
    \caption{Flowchart alur eksekusi program lexer dari input ke output}
    \label{fig:flowchart-lexer-program}
\end{figure}

Lebih lanjut, proses \texttt{lexer.tokenize()} dapat diringkas sebagai berikut. Selama belum mencapai akhir input, lexer memeriksa state saat ini, menjalankan handler state yang sesuai, dan melakukan salah satu dari tiga aksi
\begin{enumerate}
    \item \texttt{advance} ke karakter berikutnya,
    \item \texttt{emit token} jika lexeme selesai, atau
    \item \texttt{emit error} jika pola tidak valid.
\end{enumerate}
Ketika loop selesai, \texttt{handleEof()} dipanggil untuk menyelesaikan state tertunda.

\section{Struktur Direktori Program}

\begin{itemize}
    \item \textbf{\texttt{src/}}. Direktori utama yang berisi semua source code program dalam bahasa C++. Direktori ini dibagi menjadi dua modul.
          \begin{itemize}
              \item \texttt{main.cpp}. File entry point program. Bertanggung jawab untuk parse argumen command-line, membuka dan membaca file input, menginisialisasi lexer, memanggil fungsi \texttt{tokenize()}, memformat output, dan menulis hasil ke terminal serta file output (jika ada). File ini juga menangani error handling untuk file I/O dan menentukan exit code berdasarkan ada tidaknya lexical error.

              \item \texttt{lexer/}. Subdirektori yang berisi modul lexer dan token. Didalam subdirektori ini terdapat:
                    \begin{itemize}
                        \item \texttt{token.h} dan \texttt{token.cpp}. File ini berisi modul yang mendefinisikan struktur atau class \texttt{Token}, enum \texttt{TokenType} dengan 52 jenis token plus ERROR, dan fungsi-fungsi utilitas seperti konversi tipe token ke string dan pemformatan output token.

                        \item \texttt{lexer.h} dan \texttt{lexer.cpp}. File ini berisi Modul yang mengimplementasikan class atau struktur \texttt{Lexer}. Lexer mengenkapsulasi state mesin DFA, buffer pembacaan source code, posisi karakter saat ini, dan logika-logika transisi state. Fungsi-fungsi utama meliputi \texttt{tokenize()} untuk loop tokenisasi, handler-handler state (misalnya handler identifier, angka, string, operator, komentar), dan fungsi \texttt{emitToken()/emitError()} untuk mengeluarkan token.
                    \end{itemize}
          \end{itemize}

    \item \textbf{\texttt{doc/}}. Direktori dokumentasi proyek yang berisi laporan Milestone 1 dalam format LaTeX.

    \item \textbf{\texttt{test/}}. Direktori test case yang menyimpan input, expected output, dan actual output program untuk pengujian fungsionalitas lexer. Struktur di dalam direktori ini adalah sebagai berikut.
          \begin{itemize}
              \item \texttt{milestone-1/}. Subdirektori untuk test case Milestone 1 yang berisi
                    \begin{itemize}
                        \item \texttt{input-1.txt, input-2.txt, ..., input-N.txt}. File input berisi source code dalam bahasa Arion untuk di-test.
                        \item \texttt{expected-1.txt, expected-2.txt, ..., expected-N.txt}. File output yang diharapkan untuk setiap input (hasil tokenisasi yang benar).
                        \item \texttt{output-1.txt, output-2.txt, ..., output-N.txt}. File output aktual yang dihasilkan program ketika dijalankan terhadap input.
                    \end{itemize}
          \end{itemize}

    \item \textbf{\texttt{Makefile}}. File build automation yang berisi rule-rule untuk kompilasi program.

    \item \textbf{\texttt{README.md}}. File dokumentasi yang berisi deskripsi singkat program, fitur-fitur, requirement sistem, instruksi build dan run, contoh input/output, struktur direktori, informasi penulis, dan lisensi MIT.
\end{itemize}

\section{Implementasi DFA}
\subsection{Penentuan State}

Lexer ini memiliki 15 state dasar, dimulai dari state awal yaitu \texttt{START}, dan state akhir yaitu \texttt{FOUND\_EOF}. State-state lainnya merupakan state penengah yang digunakan untuk menandai identifier yang telah ditemukan.
\begin{enumerate}
    \item \texttt{START}. State awal, menunggu karakter baru untuk menentukan jenis token berikutnya.

    \item \texttt{IN\_IDENT}. Sedang membaca identifier atau keyword.

    \item \texttt{IN\_NUMBER}. Sedang membaca digit-digit konstanta integer.

    \item \texttt{DOT\_AFTER\_NUM}. Baru saja membaca titik setelah angka untuk membedakan \texttt{period} dan bilangan riil.

    \item \texttt{IN\_REAL}. Sedang membaca bagian desimal bilangan riil.

    \item \texttt{IN\_STRING}. Sedang membaca karakter di dalam tanda kutip.

    \item \texttt{SAW\_COLON}. Melihat operator \texttt{:} yang mungkin dilanjutkan menjadi \texttt{:=}.

    \item \texttt{SAW\_LESS}. Melihat operator \texttt{\textless} yang mungkin dilanjutkan menjadi \texttt{\textless\textgreater} atau \texttt{\textless=}.

    \item \texttt{SAW\_GREATER}. Melihat operator \texttt{\textgreater} yang mungkin dilanjutkan menjadi \texttt{\textgreater=}.

    \item \texttt{SAW\_EQUAL}. Melihat operator \texttt{=} yang valid hanya jika dilanjutkan menjadi \texttt{==}.

    \item \texttt{IN\_CURLY\_COMMENT}. Sedang membaca karakter di dalam komentar kurawal.

    \item \texttt{SAW\_LPAREN}. Telah membaca kurung buka biasa \texttt{(}, yang juga merupakan penanda awal komentar dengan format \texttt{(*}.

    \item \texttt{IN\_PAREN\_STAR\_COMMENT}. Memroses isi komentar pada bagian \texttt{(* ... *)}.

    \item \texttt{SAW\_STAR\_IN\_PAREN\_COMMENT}. Di akhir komentar \texttt{(*}, memroses kemunculan tanda \texttt{*} sebagai penutup akhir komentar.

    \item \texttt{FOUND\_EOF}. Menjadi penanda selesainya pembacaan stream file \textit{source code}.
\end{enumerate}

\subsection{Tabel Transisi}
State-state tersebut kami susun menjadi sebuah diagram transisi DFA. Sebelum membentuk diagram tersebut, kami menyusun tabel yang menentukan transisi dari satu state ke state yang lain berdasarkan karakter yang dibaca. Angka pada state disusun berurutan dengan state yang telah disusun pada subsection diatas.

\begin{table}[H]
    \centering
    \resizebox{\textwidth}{!}{
        \begin{tabular}{|c*{15}{|c}|}
            \hline
            \textbf{State} & \textbf{Alphabets} & \textbf{Digits} & \textbf{Space Char} & \textbf{.} & \textbf{:} & \textbf{\textless} & \textbf{\textgreater} & \textbf{=} & \textbf{\{} & \textbf{\}} & \textbf{(} & \textbf{*} & \textbf{)} & \textbf{'} & \textbf{EOF} \\ \hline
            -> q1          & q2                 & q3              & q1                  & q1         & q7         & q8                 & q9                    & q10        & q11         & q1          & q12        & q1         & q1         & q6         & q15          \\ \hline
            q2             & q2                 & q2              & q1                  & q1         & q1         & q1                 & q1                    & q1         & q1          & q1          & q1         & q1         & q1         & q1         & q15          \\ \hline
            q3             & q1                 & q3              & q1                  & q4         & q1         & q1                 & q1                    & q1         & q1          & q1          & q1         & q1         & q1         & q1         & q15          \\ \hline
            q4             & q1                 & q5              & q1                  & q1         & q1         & q1                 & q1                    & q1         & q1          & q1          & q1         & q1         & q1         & q1         & q15          \\ \hline
            q5             & q1                 & q5              & q1                  & q1         & q1         & q1                 & q1                    & q1         & q1          & q1          & q1         & q1         & q1         & q1         & q15          \\ \hline
            q6             & q6                 & q6              & q6                  & q6         & q6         & q6                 & q6                    & q6         & q6          & q6          & q6         & q6         & q6         & q1         & q15          \\ \hline
            q7             & q1                 & q1              & q1                  & q1         & q1         & q1                 & q1                    & q1         & q1          & q1          & q1         & q1         & q1         & q1         & q15          \\ \hline
            q8             & q1                 & q1              & q1                  & q1         & q1         & q1                 & q1                    & q1         & q1          & q1          & q1         & q1         & q1         & q1         & q15          \\ \hline
            q9             & q1                 & q1              & q1                  & q1         & q1         & q1                 & q1                    & q1         & q1          & q1          & q1         & q1         & q1         & q1         & q15          \\ \hline
            q10            & q1                 & q1              & q1                  & q1         & q1         & q1                 & q1                    & q1         & q1          & q1          & q1         & q1         & q1         & q1         & q15          \\ \hline
            q11            & q11                & q11             & q11                 & q11        & q11        & q11                & q11                   & q11        & q11         & q1          & q11        & q11        & q11        & q11        & q15          \\ \hline
            q12            & q1                 & q1              & q1                  & q1         & q1         & q1                 & q1                    & q1         & q1          & q1          & q1         & q13        & q1         & q1         & q15          \\ \hline
            q13            & q13                & q13             & q13                 & q13        & q13        & q13                & q13                   & q13        & q13         & q13         & q13        & q14        & q13        & q13        & q15          \\ \hline
            q14            & q13                & q13             & q13                 & q13        & q13        & q13                & q13                   & q13        & q13         & q13         & q13        & q14        & q1         & q13        & q15          \\ \hline
            * q15          & q15                & q15             & q15                 & q15        & q15        & q15                & q15                   & q15        & q15         & q15         & q15        & q15        & q15        & q15        & q15          \\ \hline
        \end{tabular}
    }
    \caption{Tabel Transisi DFA Penuh (dari q1 hingga final marker q15)}
    \label{tab:dfa-transition}
\end{table}

\subsection{Diagram Transisi DFA}
Setelah merumuskan tabel transisi state, dibentuk diagram state DFA yang meliputi
\begin{itemize}
    \item \textbf{Initial State}: q1 (\texttt{START})
    \item \textbf{Final State}: q15 (\texttt{FOUND\_EOF}).
\end{itemize}

\begin{figure}[H]
    \centering
    \includegraphics[width=\linewidth]{03-implementasi/diagram-transisi.png}
    \caption{Diagram Transisi Program (Sumber: \href{https://drive.google.com/file/d/1aSSj8bhTI08BDqSrc8pjaBD1EtSzj6KU/view?usp=sharing}{(Link State Diagram)})}
\end{figure}

\section{Implementasi Kode}
\subsection{\texttt{token.h}}

File \texttt{token.h} mendefinisikan enum \texttt{TokenType} yang merepresentasikan 52 jenis token dalam bahasa Arion, ditambah satu token ERROR untuk penanganan kesalahan leksikal. File ini juga mendefinisikan struct \texttt{Token} yang menyimpan informasi lengkap tentang setiap token yang dihasilkan lexer, serta deklarasi fungsi-fungsi utilitas untuk konversi dan pemformatan.

\begin{lstlisting}[language=C++, numbers=left, numberstyle=\tiny, caption={Snippet enum TokenType dari token.h}, label={lst:token-enum}]
enum class TokenType {
    // Literal
    INTCON,              // Konstanta integer
    REALCON,             // Konstanta bilangan riil
    CHARCON,             // Konstanta karakter
    STRING,              // Sekuens karakter
    
    // Operator
    PLUS, MINUS, TIMES,  // Aritmatika dasar
    IDIV, RDIV, IMOD,    // Pembagian dan modulo
    EQL, NEQ, GTR, GEQ,  // Perbandingan
    LSS, LEQ,
    
    // Keyword
    VARSY, PROGRAMSY, IFSY, BEGINSY, ENDSY,
    WHILESY, FORSY, // ... total 52 token
    
    // Error handling
    ERROR
};
\end{lstlisting}

\textbf{Struktur Token}. Struct \texttt{Token} adalah tipe data yang menyimpan informasi lengkap tentang token yang dihasilkan lexer.

\begin{lstlisting}[language=C++, numbers=left, numberstyle=\tiny, caption={Struktur Token dari token.h}, label={lst:token-struct}]
struct Token {
    TokenType type;      // Jenis token (dari enum TokenType)
    std::string value;   // Nilai token untuk ident, 
    //intcon, realcon, charcon, string
    int line;            // Nomor baris (1-based)
    int column;          // Nomor kolom (1-based)
};
\end{lstlisting}

\textbf{Fungsi Utilitas}. File ini juga mendeklarasikan tiga fungsi utilitas penting sebagai berikut.
\begin{itemize}
    \item \texttt{std::string tokenTypeToString(TokenType type)} . Fungsi ini digunkan untuk mengonversi enum TokenType menjadi string untuk ditampilkan pada output. Contohnya, \texttt{INTCON} dikonversi menjadi \texttt{"intcon"}.

    \item \texttt{bool tokenHasValue(TokenType type)}. Fungsi ini digunakan untuk mengecek apakah token memiliki nilai yang perlu ditampilkan. Contohnya adalah IDENT, INTCON, REALCON, CHARCON, STRING memiliki nilai, sedangkan PLUS, SEMICOLON tidak.

    \item \texttt{std::string formatToken(const Token\& token)}. Fungsi ini digunakan untuk Memformat token untuk output sesuai spesifikasi format lexer. Jika token memiliki nilai, format output adalah \texttt{<tipe-token>(<nilai>)}, misalnya \texttt{ident(x)}. Jika tidak ada nilai, hanya \texttt{<tipe-token>}, misalnya \texttt{plus}.
\end{itemize}

\subsection{\texttt{token.cpp}}

File \texttt{token.cpp} mengimplementasikan tiga fungsi pembantu yang dideklarasikan di \texttt{token.h}.

\begin{enumerate}
    \item \textbf{tokenTypeToString()}. Fungsi ini mengkonversi nilai enum \texttt{TokenType} menjadi representasi string. Implementasinya menggunakan \texttt{switch-case} untuk memetakan setiap token type ke string yang sesuai, semua dalam huruf kecil sesuai format output yang ditentukan. Contohnya, \texttt{TokenType::INTCON} dikonversi menjadi \texttt{"intcon"}, \texttt{TokenType::PROGRAMSY} menjadi \texttt{"programsy"}, dan seterusnya.
    \item \textbf{tokenHasValue()}. Fungsi ini mengecek apakah suatu tipe token memiliki nilai yang perlu ditampilkan dalam tanda kurung. Token dengan nilai antara lain IDENT (identifier), INTCON (konstanta integer), REALCON (konstanta riil), CHARCON (konstanta karakter), STRING (string literal), COMMENT (komentar), dan ERROR (pesan error). Sementara kategori token lain seperti operator, delimiter, dan keyword tidak memiliki nilai yang ditampilkan.
    \item \textbf{formatToken()}. Fungsi ini memformat token untuk output. Jika token memiliki nilai (berdasarkan \texttt{tokenHasValue()}), format outputnya adalah \texttt{<tipe> \\(<nilai>)}, misalnya \texttt{ident(Hello)} atau \texttt{intcon(42)}. Jika token tidak memiliki nilai, format outputnya hanya \texttt{<tipe>}, misalnya \texttt{plus} atau \texttt{semicolon}.
\end{enumerate}
Snippet dari \texttt{token.cpp} ditunjukkan berikut.

\begin{lstlisting}[language=C++, numbers=left, numberstyle=\tiny, caption={Implementasi formatToken() dari token.cpp}, label={lst:format-token}]
std::string formatToken(const Token& token) {
    if (tokenHasValue(token.type)) {
        return tokenTypeToString(token.type) + " (" + 
               token.value + ")";
    }
    return tokenTypeToString(token.type);
}
\end{lstlisting}

\subsection{\texttt{lexer.h}}

File \texttt{lexer.h} mendeklarasikan enum \texttt{State} yang merepresentasikan 13 state dalam DFA lexer, serta kelas \texttt{Lexer} yang mengimplementasikan mesin DFA.

\textbf{Enum State}. Enum ini mendefinisikan semua state yang ada dalam DFA sebagai berikut.

\begin{lstlisting}[language=C++, numbers=left, numberstyle=\tiny, caption={Snippet enum State dari lexer.h}, label={lst:state-enum}]
enum class State {
    START,                   // Menunggu karakter baru
    IN_IDENT,                // Sedang membaca identifier
    IN_NUMBER,               // Sedang membaca digit integer
    DOT_AFTER_NUM,           // Melihat '.' setelah digit
    IN_REAL,                 // Sedang membaca desimal float
    IN_STRING,               // Sedang membaca string literal
    SAW_COLON,               // Melihat ':'
    SAW_LESS,                // Melihat '<'
    SAW_GREATER,             // Melihat '>'
    SAW_EQUAL,               // Melihat '='
    IN_CURLY_COMMENT,        // Di dalam komentar kurawal
    SAW_LPAREN,              // Melihat '('
    IN_PAREN_STAR_COMMENT,   // Di dalam komentar '(*'
    SAW_STAR_IN_PAREN_COMMENT, // Melihat '*' di dalam komentar
    FOUND_EOF                // Penanda akhir file (EOF)
};
\end{lstlisting}

\textbf{Kelas Lexer}. Kelas ini mengenkapsulasi semua data dan logika untuk melakukan tokenisasi. Data member privat meliputi data sebagai berikut.
\begin{itemize}
    \item \texttt{source}. String berisi seluruh source code yang akan di-tokenisasi.
    \item \texttt{pos, line, col}. Posisi pembacaan saat ini (posisi karakter, nomor baris, nomor kolom).
    \item \texttt{state}. State DFA saat ini.
    \item \texttt{buffer}. Akumulasi karakter untuk token yang sedang diproses.
    \item \texttt{tokenStartLine, tokenStartCol}. Posisi awal token yang sedang diproses (untuk keperluan error reporting).
    \item \texttt{errorFlag}. Flag untuk menandai adanya error leksikal.
    \item \texttt{tokens}. Daftar (vector) hasil token yang dihasilkan.
\end{itemize}

Fungsi publik meliputi fungsi-fungsi di bawah ini.
\begin{itemize}
    \item \texttt{Lexer(const std::string\& source)}. Konstruktor yang menginisialisasi lexer dengan source code.
    \item \texttt{std::vector<Token> tokenize()}. Fungsi utama yang menjalankan DFA dan mengembalikan daftar token hasil analisis.
    \item \texttt{bool hasErrors() const}. Mengecek apakah ada error leksikal.
\end{itemize}

Fungsi privat meliputi handler untuk setiap state (\texttt{handleStart()}, \texttt{handleInIdent()}, \texttt{handleInNumber()}, dll.), serta fungsi-fungsi utilitas seperti \texttt{emitToken()}, \texttt{emitError()}, dan \texttt{lookupKeyword()}.

\subsection{\texttt{lexer.cpp}}

File \texttt{lexer.cpp} mengimplementasikan kelas \texttt{Lexer}, termasuk fungsi utilitas privat dan handler untuk setiap state DFA. Implementasi mengikuti pola \textit{state machine} klasik. Polanya adalah untuk setiap state, handler yang sesuai membaca karakter saat ini, memutuskan tindakan dan melanjutkan ke karakter berikutnya.

\textbf{Tabel Keyword}. Di awal file terdapat \texttt{static const std::map<>} yang menyimpan 27 keyword dari bahasa Arion. Perbandingan dilakukan case-insensitive dengan mengkonversi identifier ke lowercase sebelum lookup. Contohnya, identifier \texttt{PROGRAM} akan dicocokkan dengan key \texttt{program} dan menghasilkan token \texttt{PROGRAMSY}.

\textbf{Konstruktor dan Fungsi Utilitas}. Konstruktor \texttt{Lexer()} menginisialisasi semua member variable. Fungsi utilitas \texttt{current()}, \texttt{isAtEnd()}, dan \texttt{advance()} menangani pembacaan karakter dan update posisi line/column. Fungsi \texttt{emitToken()} dan \texttt{emitError()} menambahkan token/error ke daftar output. Fungsi \texttt{lookupKeyword()} mencari identifier di tabel keyword.

\textbf{Fungsi tokenize() Utama}. Fungsi ini menjalankan loop utama yang membaca hingga akhir input, memanggil handler state yang sesuai, dan setelah akhir file, memanggil \texttt{handleEof()} untuk menyelesaikan token yang belum selesai. Detail dari \texttt{tokenize()} dapat dilihat di Snippet kode di bawah ini.
\begin{lstlisting}[language=C++, numbers=left, numberstyle=\tiny, caption={Loop utama tokenize() dari lexer.cpp}, label={lst:tokenize-loop}]
std::vector<Token> Lexer::tokenize() {
    while (!isAtEnd()) {
        switch (state) {
            case State::START:     handleStart(); break;
            case State::IN_IDENT:  handleInIdent(); break;
            case State::IN_NUMBER: handleInNumber(); break;
            // ... handler untuk state lainya ...
        }
    }
    handleEof();  // Tangani token yang belum selesai
    return tokens;
}
\end{lstlisting}

\textbf{Handler State}. Setiap state memiliki handler yang bertanggung jawab untuk logika transisi DFA. Contohnya:

\begin{itemize}
    \item \textbf{handleStart()}. Handler ini Membaca karakter pertama dan menentukan jenis token. Jika huruf, transisi ke \texttt{IN\_IDENT}. Jika digit, transisi ke \texttt{IN\_NUMBER}. Jika tanda kutip, transisi ke \texttt{IN\_STRING}. Operator dan delimiter single-character langsung emit token. Whitespace diabaikan.

    \item \textbf{handleInIdent()}. Handler ini akan mengakumulasi huruf/digit ke buffer. Ketika melihat non-alphanumeric, lookup keyword dan emit token IDENT atau keyword, lalu kembali ke \texttt{START}.

    \item \textbf{handleInNumber()}. Handler ini akan mengakumulasi digit. Ketika melihat `.`, transisi ke \texttt{DOT\_AFTER\_NUM} untuk menentukan apakah itu bilangan riil atau period. Non-digit menyebabkan emit INTCON dan kembali ke \texttt{START}.

    \item \textbf{handleDotAfterNum()}. Jika digit berikutnya ada, gabungkan ke buffer dan transisi ke \texttt{IN\_REAL}. Jika bukan digit, emit INTCON, emit PERIOD, dan kembali ke \texttt{START}.

    \item \textbf{handleSawColon(), handleSawLess(), dst.}. Handler ini untuk operator multi-karakter. Misalnya, setelah \texttt{:} dilihat, jika karakter berikutnya adalah \texttt{=}, emit BECOMES; sebaliknya emit COLON.

    \item \textbf{handleInString()}. Akumulasi karakter dalam tanda kutip. Jika menemukan tanda kutip penutup, tentukan apakah string (panjang > 1) atau CHARCON (panjang = 1).

    \item \textbf{handleInCurlyComment()}. Akumulasi karakter yang berada di dalam komentar bergaya kurung kurawal \texttt{\{ \}}. Program akan terus membaca hingga menemukan kurung kurawal penutup \texttt{\}}, kemudian emit COMMENT dan kembali ke \texttt{START}.

    \item \textbf{handleSawLparen()}. Handler ini mengevaluasi karakter setelah kurung buka biasa \texttt{(}. Jika karakter berikutnya adalah asterisk \texttt{*}, ini berarti awal dari komentar multi-baris \texttt{(*}, sehingga state bertransisi ke \texttt{IN\_PAREN\_STAR\_COMMENT}. Jika bukan, program emit LPARENT dan kembali ke \texttt{START}.

    \item \textbf{handleInParenStarComment()}. Akumulasi karakter di dalam blok komentar \texttt{(* ... *)}. Jika program mendeteksi karakter \texttt{*}, state akan berubah menjadi \texttt{SAW\_STAR\_IN\_PAREN\_COMMENT} untuk menunggu hingga karakter selanjutnya adalah \texttt{)}.

    \item \textbf{handleSawStarInParenComment()}. Mengevaluasi karakter persis setelah asterisk di dalam komentar. Jika karakter tersebut adalah \texttt{)}, blok komentar dinyatakan selesai, program emit COMMENT lalu kembali ke \texttt{START}. Jika karakternya masih berupa \texttt{*}, program tetap di state ini. Karakter selain itu akan mengembalikan state ke \texttt{IN\_PAREN\_STAR\_COMMENT}.

    \item \textbf{handleEof()}. Handler ini bertugas sebagai pembersih saat program mencapai akhir file (\textit{End of File}). Ia akan memproses sisa karakter yang tertinggal di buffer, seperti mencetak identifier atau angka terakhir. Prosedur ini juga akan melemparkan (\textit{error}) jika ada string atau komentar yang terputus (belum tertutup) sebelum file benar-benar berakhir.
\end{itemize}

Contoh snippet dari \texttt{handleStart()} ditunjukkan sebagai berikut, memperlihatkan bagaimana handler menentukan transisi berdasarkan karakter saat ini.

\begin{lstlisting}[language=C++, numbers=left, numberstyle=\tiny, caption={Snippet handleStart() dari lexer.cpp}, label={lst:handle-start}]
void Lexer::handleStart() {
    char c = current();
    
    if (std::isspace(c)) {
        advance();  // Abaikan whitespace
    } else if (std::isalpha(c)) {
        buffer = c;
        state = State::IN_IDENT;  // Mulai baca identifier
        advance();
    } else if (std::isdigit(c)) {
        buffer = c;
        state = State::IN_NUMBER;  // Mulai baca bilangan
        advance();
    } else if (c == '+') {
        emitToken(TokenType::PLUS);
        advance();
    } else if (c == ':') {
        state = State::SAW_COLON;  // Bisa : atau :=
        advance();
    } else {
        // Karakter tidak dikenali
        emitError("Unexpected character");
        advance();
    }
}
\end{lstlisting}

Dengan struktur \textit{state machine} ini, lexer dapat mengenali semua pola token dalam bahasa Arion dengan cara yang terstruktur dan mudah dipahami. Setiap transisi state mewakili suatu tahap dalam proses pengenalan token, dan handler memastikan akumulasi buffer serta penentuan token type yang tepat sesuai pola yang ditemukan.